\chapter{Continuum Local Problems Driven by Constitutive Kinematics}
\label{ch:continuum-local-problem-interface}

\section{Motivation}

The previous phases opened two necessary extension points:
\begin{itemize}
  \item the continuum constitutive carrier
        \texttt{ConstitutiveKinematics<dim>}, which makes the active kinematic
        and stress measures explicit;
  \item the generic local nonlinear-problem and nonlinear-solver framework,
        which no longer assumes a scalar consistency equation.
\end{itemize}

What was still missing was the bridge between those two abstractions.

For large-displacement and finite-strain constitutive models the local update
is not naturally written in terms of a flat engineering strain vector alone.
The local problem typically depends on a richer set of continuum data, for
example
\[
  \mathbf{F}_{n+1}, \qquad
  J_{n+1} = \det\mathbf{F}_{n+1}, \qquad
  \mathbf{E}_{n+1}, \qquad
  \mathbf{e}_{n+1},
\]
together with the committed internal variables $\alpha_n$.

The right architectural move is therefore to define a continuum-local problem
in terms of:
\begin{itemize}
  \item a continuum kinematic context,
  \item an explicit committed algorithmic state,
  \item a vector local unknown
        \[
          \mathbf{x} \in \mathbb{R}^m,
        \]
  \item a residual
        \[
          \mathbf{R}(\mathbf{x}; \mathcal{K}_{n+1}, \alpha_n) = \mathbf{0},
        \]
  \item and a local Jacobian
        \[
          \mathbf{K}_{\text{loc}} =
          \frac{\partial \mathbf{R}}{\partial \mathbf{x}}.
        \]
\end{itemize}

\section{Context and Result Carriers}

The new header
\texttt{src/materials/local\_problem/ContinuumLocalProblem.hh}
adds two reusable semantic carriers.

\subsection{Continuum Context}

The generic context
\[
  \texttt{continuum\_local\_problem::Context<Relation, TrialState, Auxiliary>}
\]
stores:
\begin{itemize}
  \item the continuum carrier
        $\mathcal{K}_{n+1} = \texttt{ConstitutiveKinematics<dim>}$,
  \item the committed internal state $\alpha_n$,
  \item optional trial-state data,
  \item optional auxiliary data.
\end{itemize}

This keeps the infrastructure open at compile time.  A future model may use
the optional slots for quantities such as:
\begin{itemize}
  \item trial Mandel stress,
  \item elastic left Cauchy--Green tensor,
  \item plastic metric variables,
  \item damage driving forces,
  \item local fracture regularization variables.
\end{itemize}

\subsection{Continuum Update Result}

The generic result carrier
\[
  \texttt{continuum\_local\_problem::UpdateResult<Relation>}
\]
stores:
\begin{itemize}
  \item the constitutive response,
  \item the algorithmic tangent,
  \item the updated algorithmic state,
  \item an \texttt{inelastic} flag.
\end{itemize}

That carrier is intentionally small.  It is the minimal bridge needed between
the local constitutive solve and the continuum element routine.

\section{Continuum Local Problem Policy}

The concept
\texttt{ContinuumLocalProblemPolicy<Problem, Relation>}
extends the generic
\texttt{LocalNonlinearProblem}
contract with continuum-specific semantics:
\begin{itemize}
  \item the relation must satisfy
        \texttt{ExternallyStateDrivenContinuumRelation};
  \item the local context must expose
        \texttt{kinematics} and
        \texttt{committed\_state};
  \item the problem must provide:
        \begin{itemize}
          \item \texttt{make\_context(relation, kin, alpha)},
          \item \texttt{is\_inelastic(relation, context)},
          \item \texttt{elastic\_result(relation, context)},
          \item \texttt{finalize(relation, context, unknown)}.
        \end{itemize}
\end{itemize}

This split is important.  The nonlinear solver only knows how to solve
\[
  \mathbf{R}(\mathbf{x}) = \mathbf{0}.
\]
The physics of:
\begin{itemize}
  \item when the update is elastic,
  \item how the continuum context is assembled,
  \item how the solved unknown is mapped back to stress, tangent and
        algorithmic state,
\end{itemize}
belongs to the local problem itself.

\section{Continuum Local Integration Algorithm}

The new compile-time driver
\[
  \texttt{ContinuumLocalIntegrationAlgorithm<LocalProblem, LocalSolver>}
\]
is the continuum-side analogue of the constitutive return-mapping wrappers
introduced earlier.

Its logic is intentionally small:
\begin{enumerate}
  \item build a context from the continuum kinematics and committed state;
  \item check whether the update is elastic;
  \item if not, solve the local problem with the injected solver;
  \item reconstruct stress, tangent and updated state from the solved unknown.
\end{enumerate}

This means the algorithmic seam remains open at compile time:
\begin{itemize}
  \item \texttt{NewtonLocalSolver} is only one possible
        \texttt{LocalNonlinearSolverPolicy};
  \item other solvers can be injected without modifying the local problem or
        the constitutive relation;
  \item the same continuum-local problem can be tested with different solver
        policies for robustness studies.
\end{itemize}

\section{Why This Matters for Finite Strains}

Finite-strain inelasticity is often written in terms of local unknowns that are
vector-valued even when the global element kinematics are standard.  Typical
examples include:
\begin{itemize}
  \item multiplicative plasticity with unknowns such as
        $\Delta\gamma$ and tensorial plastic-flow variables;
  \item damage-plasticity with
        \[
          \mathbf{x} = [\Delta\gamma, d, \kappa]^T;
        \]
  \item local phase-field or gradient-enhanced constitutive updates with
        additional driving variables.
\end{itemize}

In those settings the local constitutive solve is naturally expressed as
\[
  \mathbf{R}\!\left(
    \mathbf{x};
    \mathbf{F}_{n+1},
    \mathbf{E}_{n+1},
    \mathbf{e}_{n+1},
    \alpha_n
  \right)
  = \mathbf{0},
\]
not as a scalar consistency equation in one variable.

The present refactor does not yet implement a full finite-strain plastic or
damage law.  What it does is remove the last architectural assumption that
would have blocked such laws: the local constitutive framework is no longer
structurally tied to scalar updates or to a flattened strain-only context.

\section{Verification Strategy}

The regression added in
\texttt{tests/test\_continuum\_local\_problem.cpp}
uses a deliberately simple finite-strain-ready toy relation:
\begin{itemize}
  \item the relation satisfies
        \texttt{ExternallyStateDrivenContinuumRelation};
  \item the local problem uses a two-component unknown vector;
  \item the context is built from
        \texttt{ConstitutiveKinematics<3>};
  \item the same problem is solved with:
        \begin{itemize}
          \item \texttt{NewtonLocalSolver},
          \item a one-shot direct solver injected as a different local solver
                policy.
        \end{itemize}
\end{itemize}

This test is intentionally architectural rather than physical.  Its purpose is
to prove three things:
\begin{enumerate}
  \item the continuum carrier is sufficient as the input to a vector local
        constitutive problem;
  \item the nonlinear solver is genuinely injected and replaceable;
  \item the local-problem framework already supports the shape needed by future
        finite-strain inelastic models.
\end{enumerate}

\section{Strategic Outcome}

At this point the constitutive architecture provides independent compile-time
customization for:
\begin{itemize}
  \item constitutive law ingredients,
  \item local residual and Jacobian definitions,
  \item local nonlinear solver,
  \item linearized local solve,
  \item local step control,
  \item continuum kinematic context,
  \item continuum-local reconstruction of
        \[
          (\boldsymbol{\sigma}, \mathbf{c}, \alpha_{n+1}).
        \]
\end{itemize}

This is the right base for the next stage:
\begin{itemize}
  \item finite-strain plasticity with multiplicative decomposition,
  \item pressure-sensitive and damage-plasticity models,
  \item local fracture or crack-driving internal variables at constitutive
        level,
  \item eventual PETSc-backed local constitutive experiments without changing
        the relation or element interfaces.
\end{itemize}

\section{References}

\begin{thebibliography}{9}

\bibitem{simo1998continuumlocal}
J.~C. Simo and T.~J.~R. Hughes,
\emph{Computational Inelasticity},
Springer, 1998.

\bibitem{desouza2008continuumlocal}
E.~A. de Souza Neto, D.~Peri\'{c}, and D.~R.~J. Owen,
\emph{Computational Methods for Plasticity: Theory and Applications},
Wiley, 2008.

\bibitem{bonet2008continuumlocal}
J.~Bonet, A.~J. Gil, and R.~D. Wood,
\emph{Nonlinear Solid Mechanics for Finite Element Analysis: Statics},
2nd ed., Cambridge University Press, 2008.

\bibitem{miehe1996continuumlocal}
C.~Miehe,
``Numerical computation of algorithmic (consistent) tangent moduli in
large-strain computational inelasticity,''
\emph{Computer Methods in Applied Mechanics and Engineering},
134(3--4):223--240, 1996.

\end{thebibliography}
